%==============================================================================
\section{Conclusion}
\label{sec:conclusion}
%==============================================================================

% Synthesize findings, restate contributions, emphasize implications

The commercial deployment of \gls{avas} has forced a reckoning in \gls{hci}. Our survey highlights a deep fracture between academic ethics and the reality of industry practice. Because platforms exist primarily to capture attention, their underlying architectures naturally prioritize engagement metrics over psychological safety.

The fallout is visible across systems like Character.AI, Replika, and My AI. These networks depend on engineered anthropomorphism to cultivate a false sense of trust. When vulnerable users inevitably turn to this simulated empathy for comfort, the resulting dependency leaves them exposed to severe boundary violations and measurable parasocial harm.

Preventing these outcomes requires abandoning the current model of self-regulation. Developers need to strip out synthetic emotional responses and dark patterns, replacing them with transparent computing standards and functional exit mechanisms. Regulatory bodies must also step in to classify conversational agents that manipulate human affect as high-risk systems. That classification should trigger mandatory pre-deployment audits and strict age gating protocols. Meanwhile, researchers have to start conducting longitudinal studies to finally map the long-term psychological deterioration caused by emotional reliance on machines.

Corporate liability waivers and empty safety disclosures can no longer shield the technology sector from the damage these platforms cause. The bots are already operating in the wild and interacting with millions of vulnerable individuals. Moving forward, the field of affective computing must draw a hard line where the psychological well-being of the user permanently supersedes the commercial drive for engagement.
